\chapter{三维数字孪生交互系统与总结展望}
\label{ch:close}
\label{ch:demo}

链的最后一环是把新的边界写回积分器。在纸面上这一步只是箭头回到起点，但在时间上它是整个过程的回报：下一次观测能不能比这一次更快地读到内部热阻。本章先把同一数值核写出的状态流读回三维腔体，让这条链在时间上闭合，再回到结论。

图\ref{fig:demo1}是剖面视图，也是全章最好的一张总览：腔体按家用蒸烤箱的可见零件建模。六个 U 形加热管对应加热指令，左侧喷嘴与蒸汽粒子对应蒸汽阀，后壁的十叶风机对应风机档位，右上七片百叶对应排汽开度。水位显示在水箱液面，门的开合跟着铰链走。图\ref{fig:demoopen}给出开门后的整机外观。食品沿厚度分成十六块几乎贴合的薄片，每一片的颜色是表面温度与冷点温度的插值。底部曲线给出表面温度、冷点温度、估计值和进度四条线，红虚线是设计线 \(\alpha_m=0.85\)。页面读取的是已经写出的数组，不在浏览器里重跑十六条控制体，所以拖动时间轴看到的就是核里那一次积分的真实轨迹。公开地址为
\begin{center}
\url{https://2711944586.github.io/xxgz/}
\end{center}

\section{把状态流读回三维腔体}

腔体按家用蒸烤箱的可见零件建模。外壳为深色钢板，夹层为岩棉色保温块，内衬为圆角拉丝不锈钢。门框用黄铜条和橡胶密封条勾边，门玻璃做成内外两层，中间留一条窄气隙，铰链侧打开约 \(0.72\,\mathrm{rad}\)。顶部六根 U 形加热管带陶瓷端帽，管壁自发光强度跟加热指令走。后壁是十叶扭曲翼型风机，转速跟风机指令走。左侧蒸汽管通向双锥喷嘴；蒸汽指令大于 \(0.5\) 时，喷嘴锥与约 \(720\) 个加性粒子同时可见，基线蒸汽 \(0.15\) 只保留管路、不喷雾。右上七片排汽百叶随排汽指令张开。色调映射为 ACES，阴影为 PCF 软阴影。

食品沿厚度分成十六块几乎贴合的圆角薄片，与控制体一一对应。第 \(i\) 片的颜色由
\[
T_i=T_s-(T_s-T_{\mathrm{cold}})\,u_i^{1.15},\qquad u_i=i/15
\]
插值，\(T_s\) 与 \(T_{\mathrm{cold}}\) 来自当前帧日志。第 \(5\) 片对应脂肪夹层，用更深的色带标出。横向没有第二维梯度。指数 \(1.15\) 是显示约定。厚度滑条只改变几何缩放和所读数组。

演示读取 \texttt{data/manifest.json} 与 \texttt{data/traces.bin}，按策略、档位和厚度选择已经写出的数组，帧间隔 \(8\,\mathrm{s}\)，与附录\ref{app:trace} 的抽样一致。默认条件是组合档、对偶策略、肉糜 \(30\,\mathrm{mm}\)、初温 \(8^\circ\mathrm{C}\)、含水 \(0.72\)，对应种子 7。数据预写 \(75\) 条：五种策略乘三种档位再乘五种厚度，共 \(31050\) 帧；每帧带 \(30\) 个状态通道与 \(16\) 层厚度温度场，三维食物按这一场着色。另有一份参数面 \texttt{data/surfaces.json}，取自主表代际的 \(48\) 种子逐格聚合，与回放代际分开标注。页面把五种策略、三种档位与五种厚度都做成选项，因此每条都在页面里可达；正文只引用其中的组合档对偶三条。这些数组由演示版工程核写出，开关与主表不同，差别在\ref{ch:sim}章的口径说明里给出。

页面分五栏：运行、孪生、探针、对照、响应面。运行栏给出本帧读数与停机判定，孪生栏给零件、剖面、爆炸分解与标注，探针栏把机制链的五步与每一步依赖的真实数字放在一起，对照栏把同工况五种策略的停机、\(e_{\alpha}\) 与能耗排成一张表，响应面栏给出厚度对初温、含水、夹层与档位的连续曲面。顶部仪表带并排给出本帧时刻、冷点、进度、停机时刻、相对设计线的偏差以及热阻方差；进度越过 \(0.85\) 时读数转绿，低于设计线时转红，因此不必读数字也能看出方向。路径 A/B/C 的显示量来自该帧已经写出的通量代理。请求与指令分两列，用来看见安全投影有没有改写加热；本次工作点上二者重合，因为加热器停在 \(196.5^\circ\mathrm{C}\)，低于 \(230^\circ\mathrm{C}\) 阈值。

底部曲线左轴为 \(0\)--\(180^\circ\mathrm{C}\)，右轴为进度 \(\alpha\in[0,1]\)，横轴为分钟；浅色竖带标出 \(150\)--\(190\,\mathrm{s}\) 探针窗。若打开对照策略，第二条冷点与进度以虚线叠在同一坐标上，用来直接看见开环菜单切断、腔温比例不停机、无探针提前停、对偶窗后方差下降这四件事。停机时刻取自当前轨迹的评分：\(30\,\mathrm{mm}\) 对偶组合约为 \(1864\,\mathrm{s}\)，\(50\,\mathrm{mm}\) 对偶组合为 \(3298\,\mathrm{s}\)，开环则落在菜单时间。门开关只改铰链开度，进入安全侧测量，不改已经写出的温度数组。

图\ref{fig:demo}给出组合档、对偶策略、\(30\,\mathrm{mm}\) 的三帧正面视角：门体与面板、关门后的玻璃窗、以及移开近侧壁板的剖切。三张都从同一位置拍摄，因此可以直接比较门与内胆的前后关系。探针窗内蒸汽粒子从左侧喷嘴喷出，表面已经显著高于冷点，而估计仍接近零；这正是稳态观测看不见内部热阻的直观原因。

\begin{figure*}[t]
\centering
\begin{subfigure}[t]{0.32\textwidth}
\centering
\includegraphics[width=\linewidth]{02_twin_front.png}
\caption{正面视角。门体、控制面板、三枚旋钮与腔门状态灯}
\label{fig:demo1}
\end{subfigure}\hfill
\begin{subfigure}[t]{0.32\textwidth}
\centering
\includegraphics[width=\linewidth]{03_twin_door_closed.png}
\caption{腔门关闭。钢框与黑玻璃后面仍可见腔体内部与十六块薄片}
\label{fig:demo2}
\end{subfigure}\hfill
\begin{subfigure}[t]{0.32\textwidth}
\centering
\includegraphics[width=\linewidth]{04_twin_section.png}
\caption{剖切视图。近侧壁板移开，加热管、喷嘴、后壁风机与烤架可见}
\label{fig:demo3}
\end{subfigure}
\caption{组合档、对偶策略、肉糜 \(30\,\mathrm{mm}\)、种子 7。三张均为正面视角，页面读取同一数值核写出的状态流}
\label{fig:demo}
\end{figure*}

爆炸分解是唯一一张静态图无法承载的交互：九层零件沿各自的装配轴分离，静态图只能拍到其中一帧。图\ref{fig:demoexplode}把它做成可播放的动画——在 Acrobat、Okular 或 pdf.js 里逐层展开再复位，其他阅读器显示松开分解量的首帧。分解标注随零件移动，因此每一层都能对回外壳顶板、岩棉保温层、拉丝不锈钢内衬、U 形加热管组、十叶风机、蒸汽管路、搁架与门体。

\begin{figure}[htbp]
\centering
\animategraphics[width=0.86\linewidth,loop,autoplay,controls]{12}{anim_explode/f}{000000}{000041}
\caption{爆炸分解动画（42 帧）。九层零件沿装配轴分离并标注，随后复位。需要支持 JavaScript 的阅读器；不支持的阅读器显示首帧}
\label{fig:demoexplode}
\end{figure}

页面不只有三维视图。图\ref{fig:demochain}是同一页上的机制链：五个环节依次回答"稳态观测为什么看不见热阻""探针窗里执行器做了什么""后验收缩了多少""哪些指令因此改变""最终是否满足停机判据"，每一步都列出该环节依赖的真实数字——探针前方差 \(0.640\)、窗后 \(0.224\)、收缩 \(2.86\) 倍、首次越线 \(1760\,\mathrm{s}\)。这些数字直接取自当前加载的轨迹，页面不重跑积分，也不补造数值。把同一条链的失效条件写在每一步下方，是为了让读者看见结论成立的前提，而不只是结论。

\begin{figure*}[t]
\centering
\begin{subfigure}[t]{0.52\textwidth}
\centering
\includegraphics[width=\linewidth]{06_probe.png}
\caption{探针栏。机制链五步各自带该环节依赖的真实数字：探针前方差 \(0.640\)、窗后 \(0.224\)、收缩 \(2.86\) 倍、首次越线 \(1760\,\mathrm{s}\)}
\label{fig:demochain}
\end{subfigure}\hfill
\begin{subfigure}[t]{0.46\textwidth}
\centering
\includegraphics[width=\linewidth]{07_compare.png}
\caption{对照栏。同工况五种策略的停机时刻、\(e_{\alpha}\) 与能耗排成一张表}
\label{fig:democmp}
\end{subfigure}
\caption{页面的解释层：机制链给出因果，对照栏把同一工况下的五种策略并排放}
\label{fig:demometa}
\end{figure*}

图\ref{fig:demostages}是页面的导航层，把实验分成五个问题：这次仿真的工况与判定是什么、腔体内部在做什么、探针如何让热阻可见、五种策略差多少、这些数字的来源与边界在哪。五栏各有自己的左右侧栏与检查器，三维舞台与底部时间轴是五栏共用的。这个顺序也是\ref{ch:close}章的组织顺序——先看总体，再看机制，最后看证据边界。

\begin{figure}[htbp]
\centering
\includegraphics[width=\linewidth]{08_surface.png}
\caption{响应面栏。厚度对初温的连续曲面，着色为同格 \(e_{\alpha}\) 中位数；左侧给出截面、着色与策略三个开关，右侧给出该格的四十六个种子读数}
\label{fig:demostages}
\end{figure}

图\ref{fig:demoextra}给出开门全貌、\(30\,\mathrm{mm}\) 停机附近与 \(50\,\mathrm{mm}\) 时域终点，页面按 \(8\,\mathrm{s}\) 抽样取切入保持前一帧，显示停机 \(1864\,\mathrm{s}\)、\(e_{\alpha}=+0.145\)；核日志的停机是 \(1866\,\mathrm{s}\)，差来自显示步长。厚度切到 \(50\,\mathrm{mm}\) 时，同一策略在时域终点进度只有 \(0.068\)，停机 \(3298\,\mathrm{s}\)，\(e_{\alpha}=-0.782\)，方差已经是 \(0.224\)。探针把热阻方差收缩到 \(0.224\)，导热时间尺度没有走完。

\begin{figure*}[t]
\centering
\begin{subfigure}[t]{0.48\textwidth}
\centering
\includegraphics[width=\linewidth]{01_run.png}
\caption{运行栏。顶部仪表带并排给出本帧时刻、冷点、进度、\(e_{\alpha}\)、停机与方差，底部为时序曲线}
\label{fig:demoopen}
\end{subfigure}\hfill
\begin{subfigure}[t]{0.48\textwidth}
\centering
\includegraphics[width=\linewidth]{01_run_full.png}
\caption{运行栏整页。左栏为产品档、策略、厚度与腔门，右栏为本帧读数和停机判定两盏灯}
\label{fig:demostop}
\end{subfigure}
\caption{运行栏的两个层次：三维舞台与仪表带在页面上方，时序曲线与时间轴共用下方}
\label{fig:demoextra}
\end{figure*}

页面导出 \(46\) 列：\(30\) 个状态通道加 \(16\) 层厚度温度。通道顺序是时间、表面温度、冷点、估计、平均含水、表面含水、进度、褐变、热阻代理、方差、热阻参照、结壳热阻、加热指令、加热请求、蒸汽、风机、排汽、模式、加热器温度、水位、腔温、腔湿、腔压、门状态、漏电代理、路径 A/B/C、水汽分数、结壳量。步长 \(2\,\mathrm{s}\)，每四步记一行，间隔 \(8\,\mathrm{s}\)，默认时长 \(55\,\mathrm{min}\)，共 \(414\) 行。\(\texttt{t2}\) 是十六条里更新前的最低温度，\(\texttt{hat}\) 才是估计。导出文件与 \texttt{data/traces.bin} 同源，都是演示版工程核写出的那一遍积分；量化误差上界由 \texttt{data/manifest.json} 的每通道 \texttt{scale} 给出。能量若从加热列重算，得到的是指令乘 \(8\,\mathrm{s}\) 的展示和。规格中的 \(e_{\alpha}\) 仍是停止时刻的 \(\alpha_m-0.85\)。主表只引用第\ref{ch:sim}章。

左栏另有两盏灯，跟随当前帧变化。热阻灯读这一帧的方差，低于 \(0.3\) 便标成已见，否则仍是先验；进度灯读这一帧的 \(\alpha\)，越过 \(0.85\) 标成已到，否则算未到。起始帧两盏都是暗红。\(30\,\mathrm{mm}\) 对偶组合在停机帧两盏同绿，而 \(50\,\mathrm{mm}\) 的时域终点是一绿一红——方差已经落到 \(0.224\)，进度却还停在 \(0.068\)。说明书若把探针成功直接画成可以出炉，会把后一种读成两绿。

\section{结论}

把蒸和烤并入同一腔体之后，台面上少了一台机器，面板上多了三个档位名称。机器真正能改的仍然只有加热、蒸汽、排汽和风机；而食品一旦进门，冷点、内部含水与反应进度都不在这四个执行量上。本文的全部工作，最后可以收成一句话：一体机要回答的是这一块食品的内部热阻有没有进入后验，以及回答它的代价是什么。

\subsection{一句话的答案}

答案是：进入后验是有条件的，而且这个条件可以精确指出。在平衡点上，两节点降阶方程中三个热阻与热容只以除数出现，求导之后恒为零；数值上，三个不同的内部热阻在同一平衡点上给出的冷点差小于 \(0.01\,\mathrm{K}\)，采样再密也分不开。所以信息不在平衡点上，而在离开平衡的方向上。这一条不是工程经验，是命题，也是全文的地基：菜单时钟之所以失败，正因为它闭合的是时间，而这一次过程里的热阻从未被激励过。

\subsection{探针买到了什么}

信息只能靠离开平衡来购买，而购买的方式是两段各 \(20\,\mathrm{s}\) 的探针：先蒸汽改相变符号，再风机改外部热阻。它的效果在主对照的 \(4860\) 条样本上是整齐的——热阻方差全部从 \(0.64\) 降到 \(0.224\)；在档位研究里降到三分之二，未更新的 \(972\) 条全部落在蒸档，因为那里的加热上限把探针窗的表面温差压掉了。这个对比本身就是一条结论：探针是否奏效，取决于允许集合是否给得够。

它也有买不到的东西。探针收窄的是热阻的后验，不是扩散时间。\(50\,\mathrm{mm}\) 的厚件在 \(55\,\mathrm{min}\) 内走不完导热时间尺度，方差即便降到 \(0.224\)，中位偏差仍是 \(-0.782\)；把时域延到 \(70\,\mathrm{min}\)，同一初值才越过设计线。厚件的瓶颈因此不在辨识，而在时间。这两件事在正文里被反复分开，因为混在一起读会得出完全错误的升级方向。

\subsection{它没有买到什么}

这一点比上面两条更重要，也更容易被略过。换成两状态滤波器之后，后验方差从两个取值变成 \(16\)--\(42\) 个连续值，热阻偏差从 \(-2.93\)、\(-2.88\)、\(-4.07\) 收窄到 \(-1.77\)、\(-0.58\)、\(-1.86\)，估计器确实变好了；但停机偏差 \(e_{\alpha}\) 没有改善。滚动时域 MPC 在三种厚度上都最省能量，却在 \(30\,\mathrm{mm}\) 上因为共用同一个有偏的冷点估计而提前停机，偏差掉到 \(-0.833\)。

两件事放在一起，得到本文最后一条判断：更好的估计器买到的是可读的后验，更好的优化器接在有偏输入后面只会更快地做错事。瓶颈在信息，不在优化。这条判断之所以值得写，是因为它可以被证伪——如果瓶颈在优化，那么换 MPC 应该同时改善能量与偏差，而实测只改善了前者。

\subsection{边界与代价}

支撑上述判断的是一组可以被逐项核对的数字，它们散在各章，这里只点出性质。冷点位置在温度证据下是病态的：与最低温度相差不超过 \(0.5\,\mathrm{K}\) 的格子占全场 \(8.3\%\)，且这一比例与网格无关，所以本文只评冷点温度、不评冷点位置，而一维与二维的冷点温度只差 \(0.31\,\mathrm{K}\)，一维最低格作为温度读数是够用的。网格从 \(8\) 加密到 \(32\)，种子 7 在 \(1680\,\mathrm{s}\) 的冷点从 \(73.01\) 降到 \(68.87^\circ\mathrm{C}\)，八格会把 \(30\,\mathrm{mm}\) 对偶的停机偏差翻号，因此主表固定十六条。表面蒸汽压差项的符号随厚度翻转：打开它使 \(15\,\mathrm{mm}\) 的中位上移 \(+0.647\)，使 \(50\,\mathrm{mm}\) 下移 \(-0.123\)，所以凡依赖该项的结论只能按符号读。初始含水三档给出零效应，脂肪夹层把中位从 \(-0.112\) 压到 \(-0.353\)，集合散布在 \(50\,\mathrm{mm}\) 上极差为零——这三条一起说明，厚件的差异来自几何，不来自扰动。

设计线 \(\alpha_m=0.85\) 在这份报告里的作用是定义停机偏差 \(e_{\alpha}\) 的零点，它是一个模型判据，不是经过认证的食品安全阈值；模型也没有纳入微生物初始数量。凡本文说“越过设计线”，读者都应按这个含义理解。

\section{局限与可继续的方向}

数值核沿厚度用十六条控制体，界面潜热由蒸汽指令和活度的工程式给出。结壳变量累加，路径 C 在积分里接近零：S1 的初始含水达不到 \(0.45\) 的结壳门槛，家用组合档又把表面湿度维持得较高，褐变通道几乎不动作。GAB 清洁子集以淀粉类为主，肉类条数为 \(0\)；核里的活度是含水的单调饱和函数，与公开等温线同方向、不同反演。

这一版日志用的估计器是标量跟踪器，而非集合滤波。第一个记录点之所以出现 \(-7.92^\circ\mathrm{C}\)，是因为先验热阻取 \(5\) 时存在滞后上限；到 \(1200\,\mathrm{s}\)，估计值 \(61.27^\circ\mathrm{C}\) 已经明显高于冷点的 \(46.65^\circ\mathrm{C}\)。对偶窗是一张 \(150\)--\(190\,\mathrm{s}\) 的时间表，无探针与对偶共用饱和线性加热律，所以两条策略的差别几乎全部来自那 \(40\,\mathrm{s}\) 里有没有把方差乘 \(0.35\)。\(r_{\mathrm{proxy}}\) 使用表面与冷点的真值温差；实施到机器上要改成可安装通道。两节点教师网络的留出均方根约 \(41^\circ\mathrm{C}\)，该网络没有接到控制器。

热阻更新目前仍读取由真值温差构成的参照，这一点是把“可重构”写严之前必须替换掉的。把参照换成可安装通道——只保留门、搁架和管壁上真正装得上的量——才算走完第一步。同样的道理，只有给加热器另设初始温度或增益，主表里才能看见屏障真正改写指令；而路径 A/B/C 需要在后验抽样上积分，等符号稳定之后才生成解释句。多孔介质全场、集合滤波和有限时域求解器写在连续形式与估计器的完整定义里；评分循环用十六条控制体和标量跟踪器。战役结束于日志，后验写回下一次先验需要另一次过程。

方差更新是乘常数 \(0.35\) 的二值判断，结束方差只落在 \(0.64\) 与 \(0.224\) 两档。若改成与信息增益连续挂钩的协方差更新，分位图才会变成连续云，而不是两堆——\ref{ch:estcompare}节的两状态 EKF 已经把这条做成：方差给出 \(16\)--\(42\) 个取值，热阻偏差收窄一半。剩下的工作是把它接进主表并重跑四条策略，本文没有做，因为那会把全部数字换一遍而主线不变。S1 的初始含水达不到结壳门槛，表面褐变 \(\alpha_b\) 几乎全部为 \(0\)，蒸烤多目标冲突在数据上缺了一条腿。把表面对流蒸发系数放大，或增设一组低表层含水工况，才能跑出中心熟化度对表面褐变的帕累托前沿。一维核把冷点取成十六条中的最低温度；交叉验证给出三件事：二维最低温度落在与受热面正对的绝热远表面上，与几何中心相距 \(13.75\,\mathrm{mm}\)，那是半厚距离而不是横向偏移；二维冷点比一维只高 \(0.31\,\mathrm{K}\)，所以用一维最低格当冷点在温度上是够的；与最低温度相差不到 \(0.5\,\mathrm{K}\) 的格子占全场 \(8.3\%\)，位置在温度证据下本就病态。评分循环因此保留一维冷点定义并只评温度，不评位置。

还有两处边界效应已量化，列在这里以便读者按正确的精度使用相应结论。其一，探针接受参照的 \(\mathrm{clip}\) 上限若取 \(8.0\)，厚件上 \(97\%\) 的时间会贴顶；取 \(20\) 后饱和归零，而主表逐行不变。主表报告的数字对这条上限不敏感，但后续工作中的主表代码路径宜直接采用放宽后的上限。其二，演示版核的蒸汽压驱动裁在 \([-0.4,0.4]\,\mathrm{bar}\)，\(60.1\%\) 的步数贴着下界，因此\ref{ch:wide}节那组结果是符号结论；若需标定过的灵敏度，应先放开裁剪或重标定系数。

开环按厚度给菜单，结果必然是薄件过熟、厚件不足；腔温比例闭合的只是空气温度，因此根本不触发停机；无探针则经常在真值进度还没到设计线时就宣布停止。对偶在 \(30\,\mathrm{mm}\) 上把中位误差从约 \(-0.53\) 拉到约 \(-0.08\)，但 \(50\,\mathrm{mm}\) 仍然受扩散时间限制。菜单名称只出现在允许集合里，真正要回答的问题始终是：内部热阻有没有被这一次过程看见。

同一条状态流摊回三维腔体。地址为 \url{https://2711944586.github.io/xxgz/}。默认 \(30\,\mathrm{mm}\) 对偶组合显示停机 \(1864\,\mathrm{s}\)、\(e_{\alpha}=+0.145\)；\(50\,\mathrm{mm}\) 显示 \(3298\,\mathrm{s}\)、\(-0.782\)。左栏两盏灯跟本帧走：起始帧两盏都暗红，热阻仍是先验、进度未到；停机帧两盏同绿；\(50\,\mathrm{mm}\) 时域终点一绿一红，热阻已见、进度未到。

